% ============================================================
%  ch1_intro  一、引言（原书 pp. 2229--2230）
% ============================================================

\section{引言}

\origpage{2229}1987 年，Hitchin \cite{hitchin} 考察 4 维 Yang--Mills 方程在二维上的约化，引入并研究了黎曼曲面上某一类自对偶方程的解所构成的模空间。结果发现，这个模空间具有“极其丰富的几何结构”。

设复约化结构群为 $\G$（Hitchin 最初只考虑 $\G=\SL_2$）。解模空间 $\M$ 上带有一个典范的超 Kähler 度量 $g$，其复结构记作 $I$、$J$、$K$，相应的 Kähler 形式记作 $\omega_I$、$\omega_J$、$\omega_K$。在复结构 $I$ 下，$\M$ 与 Higgs 丛——或称 Hitchin 对——$(E,\Phi)$ 的模空间 $\M_\Dol$ 一致：这里 $E$ 是 $\G$-丛，\textbf{Higgs 场} $\Phi\in H^0(C;\ad(E)\otimes K)$ 是伴随丛扭上复曲线 $C$ 上典范丛后的一个截面，于是有
\[
(\M,I)\cong \M_{\Dol}.
\]
在把 Higgs 场乘以 $i$ 所诱导的同构之下，Kähler 流形 $(\M,J,\omega_J)$ 与 $(\M,K,\omega_K)$ 同构；而它们又都同构于 $C$ 上平坦 $\G$-联络的模空间 $\M_\DR$：
\begin{equation}\label{JK}
(\M,J)\cong (\M,K)\cong \M_{\DR}.
\end{equation}
记号 $\M_\Dol$ 取 Dolbeault 之意，$\M_\DR$ 取 De Rham 之意，这一用法沿自 Simpson \cite{simpson}：正是他引入了非交换 Hodge 理论的观点来研究 $\M$。

Hitchin \cite{hitchin-stable} 随后引入 \textbf{Hitchin 映射}
\[
h:\M\to \A.
\]
当 $\G=\GL_n$ 时，它就是 Higgs 场的特征多项式
\[
h(E,\Phi)=\det(x-\Phi)\in \calA:=\bigoplus_{i=1}^n H^0(C;K^i).
\]
对一般的 $\G$，需要取 $\fg:=\mathrm{Lie}(\G)$ 上的不变多项式，再在 Higgs 场上求值。Hitchin 在 \cite{hitchin-stable} 中证明：关于 $I$-全纯辛形式 $\omega_\C:=\omega_J+i\omega_K$，$h$ 构成一个代数完全可积的 Hamilton 系统。因此 $h$ 有时被称为 \textbf{Hitchin 系统}。这意味着 $\dim(\A)=\dim(\M)/2$，并且 $h$ 的各分量函数彼此独立、Poisson 交换。此外，Hitchin 映射还是 proper 的：Hitchin 在 \cite{hitchin-stable} 中对 $\SL_2$ 证明，Nitsure \cite{nitsure} 与 Simpson \cite{simpson} 对 $\GL_n$ 证明，Faltings \cite{faltings} 对一般的 $\G$ 证明。完全可积性与 proper 性合起来推出：$h$ 的一般纤维是阿贝尔簇上的主齐性空间。特别地，它们在拓扑上同构于紧环面。

由于构造上的灵活性（曲线与结构群可以变，分支数据的类型也可以变），Hitchin 系统已与已知的大部分可积系统建立了联系 \cite{donagi-markman,donagi}。因此，它们在可积系统领域居于中心地位。

本篇综述的主要兴趣在于：镜像对称与 Langlands 对偶如何同 Hitchin 系统相关联。2003 年的论文 \cite{hausel-thaddeus} 把 Langlands 对偶群的 Hitchin 系统与镜像对称在数学上联系了起来，特别地，它针对某些 $\SL_n$ 与 $\PGL_n$ 的 Hitchin 系统提出了一条拓扑镜像对称猜想。

2007 年，Kapustin--Witten \cite{kapustin-witten} 把 Hitchin 系统放进某个超对称 $4$ 维 Yang--Mills 理论约化到 $2$ 维的框架之中。这项工作还把镜像对称与几何 Langlands 纲领理解为 $4$ 维 $S$ 对偶的约化，并对此给出了细致的刻画。由此引出了许多论文，例如 \cite{baraglia-schaposnik,garcia-prada-etal,hitchin-characteristic,hausel-mellit-pei,hausel-hitchin}，讨论 Langlands 对偶的 Hitchin 系统中成对的镜像膜。其中最后两篇强调了 $\M$ 上的一个进一步结构，即由 $(E,\Phi)\mapsto (E,\lambda\Phi)$ 给出的典范 $\T$-作用。

2010 年，Ng\^{o} \cite{ngo} 通过对 Hitchin 映射某些奇异纤维的上同调作细致分析，证明了 Langlands 纲领中的基本引理。在 \cite{hausel-gths} 中，作者讨论了 \cite{hausel-thaddeus} 的拓扑镜像对称猜想与 Ng\^{o} 这类上同调结果之间的相似之处，并猜想二者彼此相关。在 \cite{groechenig-etal1} 中，Gr\"ochenig--Wyss--Ziegler 用算术 $p$ 进积分技术证明了 \cite{hausel-thaddeus} 的拓扑镜像对称猜想；随后在 \cite{groechenig-etal2} 中，这三位作者又用这些新的 $p$ 进技术重新证明了 Ng\^{o} 的上同调结果。新近，Maulik--Shen \cite{maulik-shen} 完成了 \cite{hausel-gths} 中的部分设想，并由 Ng\^{o} 的结果推出了拓扑镜像对称猜想的一个证明。

下面我们先讨论这些发展的一些背景，然后在 \S\ref{enhanced} 中解释若干未发表的结果，内容涉及 Langlands 对偶群在幂零锥尖端处的增强镜像对称。
